\SourceLeaf{24}
\[
[6]\qquad X=\frac{4a\pi BC}{AA}\int\frac{tt\mathrm{d}t}{\BareRoot{((1-(1-\frac{BB}{AA})tt)(1-(1-\frac{CC}{AA})tt)}}
\]
integratione a $t=0$ usque ad $t=1$ extensa. Manifesto attractiones axibus coordinatarum $y,z$ parallelae hinc sponte derivantur, si $a,A$ cum $b,B$ vel cum $c,C$ permutantur.

Haec itaque formula suppeditat attractionem omnium punctorum intra sphaeroidem, et quum rigorose sit vera, quantumvis proximum sit punctum attractum ipsi sphaeroidis superficiei, etiam usque ad puncta in superficie posita valebit. Ad quam quum attractio punctorum externorum iam reducta sit, problema iam complete est solutum.

Aequatio $[6]$ praeterea docet, pro puncto interno attractionem omnium sphaeroidum similium similiterque positarum prorsus identicam esse. Quodsi itaque huiusmodi sphaerois in plura strata divisa concipiatur, quorum superficies internae et externae superficiei sphaeroidis sint similes similiterque positae, manifesto singula strata punctum attractum circumvolventia ad attractionem in hoc punctum nihil conferent, ita ut tantummodo restet attractio nuclei interioris. cuius superficies per ipsum punctum transit.

\Article{14}
\par
De ipsa integratione formulae $[6]$ non opus est prolixe disserere. Constat scilicet, eam a transcendentibus pendere, circulo logarithmisque altioribus, si omnes $A,B,C$ sint inaequales: in hoc itaque casu ad series confugiemus, quae tanto citius convergent, quo minus sphaerois a sphaera discrepat. Si vero duae quantitatum $A,B,C$ sunt aequales,, e. g. $A=B$, in quo casu sphaerois orta erit per revolutionem circa axem $=2C$, erit
\[
\begin{aligned}
X&=\frac{4\pi aC}{A}\int\frac{tt\mathrm{d}t}{\BareRoot{(1-(1-\frac{CC}{AA})tt)}}\\
&=\frac{2\pi a\cos\varphi}{\sin\varphi^3}(\varphi-\tfrac12\sin2\varphi)
\end{aligned}
\]
statuendo $\frac{C}{A}=\cos\varphi$, vel $\BareRoot{(1-\frac{CC}{AA})}=\sin\varphi$, si $C<A$, aut
\[
X=\frac{2\pi aCC}{CC-AA}-\frac{2\pi aAAC}{(CC-AA)^{\frac32}}\log\frac{C+\BareRoot{(CC-AA)}}{A}
\]
si $C>A$.
